% =====================================================================
%  Wspólna preambuła szablonów pism ZUS (świadczenia rodzicielskie 2026)
%  Silnik: XeLaTeX albo Tectonic (pdflatex NIE obsłuży polskich znaków w fontspec).
%  Kompilacja:  tectonic -X compile plik.tex
% =====================================================================
\documentclass[11pt,a4paper]{article}

\usepackage{fontspec}
\usepackage{polyglossia}
\setmainlanguage{polish}
\setotherlanguage{english}
% font domyślny (Latin Modern) — Tectonic go dostarcza

\usepackage[margin=2.5cm,a4paper]{geometry}
\usepackage{graphicx}
\usepackage{xcolor}
\usepackage{array}
\usepackage{enumitem}
\usepackage{fancyhdr}
\usepackage{hyperref}

\definecolor{zus}{HTML}{007A33}      % zieleń ZUS
\definecolor{szary}{HTML}{6B7280}
\definecolor{ramka}{HTML}{9CA3AF}

\hypersetup{colorlinks=true, linkcolor=zus, urlcolor=zus, pdftitle={Szablony pism ZUS}}

% --- pole do wypełnienia odręcznego -----------------------------------
% \Pole{nazwa} rysuje kropkowaną linię z podpisem pola pod spodem.
% Pole do wypełnienia: kropkowana linia, pod nią drobny podpis pola.
\newcommand{\Pole}[1]{%
  \textcolor{szary}{\dotfill}\,%
  {\scriptsize\textcolor{szary}{\textsc{#1}}}%
}

% --- blok warunkowy: treść pokazywana tylko, gdy pole prawdziwe --------
\newenvironment{warunek}{%
  \par\smallskip\noindent\begin{minipage}{\linewidth}%
  \color{black}%
}{%
  \end{minipage}\par\smallskip%
}

% --- nagłówek pisma ---------------------------------------------------
\newcommand{\NaglowekPisma}[1]{%
  \par\vspace{0.8em}%
  \begin{center}%
    {\large\bfseries\MakeUppercase{#1}}%
  \end{center}%
  \vspace{0.4em}%
}

% --- blok nadawcy (pierwsze linie pisma) -------------------------------
\newenvironment{nadawca}{%
  \begin{flushleft}\small
}{%
  \end{flushleft}\normalsize
}

% --- stopka z numerem strony -----------------------------------------
\pagestyle{fancy}
\fancyhf{}
\renewcommand{\headrulewidth}{0pt}
\fancyfoot[C]{\textcolor{szary}{\small\thepage}}

% --- metadane szablonu na początku pisma ------------------------------
\newcommand{\Metadane}[3]{%
  \begin{flushright}\footnotesize\textcolor{szary}%
    \textbf{#1}\\[2pt] {\itshape #2}\\[2pt] #3%
  \end{flushright}\vspace{0.5em}%
}

\setlength{\parindent}{0pt}
\setlength{\parskip}{0.6em}


\begin{document}


\begin{center}
{\LARGE\bfseries Szablony pism ZUS}\\[4pt]
{\large Świadczenia rodzicielskie --- stan prawny 2026}\\[2pt]
{\small\color{szary}16 szablonów, źródło: 89 spraw wzorcowych}
\end{center}


\tableofcontents
\clearpage


\section{Wniosek o udzielenie urlopu macierzyńskiego}

{\small\color{szary}Szablon P01 \textbullet{} art. 180 \S{} 1 k.p.}

\Pole{imie}\\[-2pt]
\Pole{adres}\\[-2pt]
PESEL: \Pole{pesel}

\Pole{platnik nazwa}
\Pole{platnik adres}
NIP: \Pole{platnik nip}

\Pole{data pisma}, \Pole{adres miejscowosc}

\NaglowekPisma{WNIOSEK O UDZIELENIE URLOPU MACIERZYŃSKIEGO}

Zgodnie z art. 180 \S{} 1 Kodeksu pracy wnoszę o udzielenie mi urlopu
macierzyńskiego w wymiarze \Pole{tygodnie} tygodni, przysługującego w związku
z urodzeniem \Pole{liczba dzieci pl} dziecka (dzieci) przy jednym porodzie.

Poród nastąpił dnia \Pole{data porodu}.
Urlop macierzyński ma trwać od dnia \Pole{urlop od} do dnia \Pole{urlop do}.

\begin{warunek}
Dziecko: \Pole{dziecko imie}, PESEL \Pole{dziecko pesel}.
\end{warunek}

\begin{warunek}
Do wniosku załączam zaświadczenie szpitalne, że dziecko urodziło się martwe.
\end{warunek}

\begin{warunek}
Jednocześnie, na podstawie art. 179\textsuperscript{1} Kodeksu pracy, wnoszę o udzielenie mi\\[-2pt]
bezpośrednio po urlopie macierzyńskim urlopu rodzicielskiego w pełnym wymiarze.\\[-2pt]
Oświadczam, że zasiłek macierzyński za cały okres będzie wynosił 81,5\% podstawy.
\end{warunek}

\Pole{imie}
......................................
(podpis pracownika)

\vfill
{\footnotesize\color{szary}\itshape Przed porodem wpisuje się przewidywaną datę porodu i uzupełnia po porodzie. Urlopy macierzyński i rodzicielski łącznie liczą się do limitu art. 180 \S{} 1 k.p. }

\clearpage

\section{Wniosek o część urlopu macierzyńskiego przed przewidywaną datą porodu}

{\small\color{szary}Szablon P02 \textbullet{} art. 180 \S{} 2 k.p.}

WNIOSEK O UDZIELENIE CZĘŚCI URLOPU MACIERZYŃSKIEGO\\[-2pt]
PRZED PRZEWIDYWANĄ DATĄ PORODU

Na podstawie art. 180 \S{} 2 Kodeksu pracy proszę o udzielenie mi części urlopu
macierzyńskiego przypadającej przed przewidywaną datą porodu, od dnia \Pole{urlop od}.

Przewidywana data porodu: \Pole{pdp}.
Wykorzystuję \Pole{tygodnie} tygodni z puli urlopu macierzyńskiego.

\Pole{imie}, \Pole{data pisma}
.....................................
(podpis)

\vfill
{\footnotesize\color{szary}\itshape Przed porodem można wykorzystać maksymalnie 6 tygodni urlopu macierzyńskiego.}

\clearpage

\section{Wniosek o zasiłek macierzyński za okres urlopu rodzicielskiego (formularz ZUR)}

{\small\color{szary}Szablon P03 \textbullet{} art. 182\textsuperscript{1}d k.p.; art. 29--31 ustawy zasiłkowej}

[FORMULARZ ZUS --- ZUR / ZAM]\\[-2pt]
Wniosek o zasiłek macierzyński za okres urlopu rodzicielskiego

Wypełnij ten wniosek, jeśli ubiegasz się o wypłatę zasiłku macierzyńskiego
za okres urlopu rodzicielskiego albo urlopu na warunkach urlopu macierzyńskiego.

ZAKRES WNIOSKU
Proszę o wypłatę zasiłku macierzyńskiego za okres od \Pole{urlop od} do \Pole{urlop do}.

Wnioskuję o zasiłek macierzyński za okres urlopu rodzicielskiego z powodu:
[ ] urodzenia się dziecka/dzieci
[ ] przyjęcia na wychowanie dziecka/dzieci w celu przysposobienia
[ ] przyjęcia na wychowanie dziecka/dzieci w ramach rodziny zastępczej
[ ] przerwania przez matkę dziecka okresu wypłaty zasiłku macierzyńskiego
[ ] posiadania przez matkę orzeczenia o niezdolności do samodzielnej egzystencji
[ ] porzucenia dziecka przez matkę
[ ] śmierci matki dziecka/dzieci
[ ] podjęcia przez matkę pracy w wymiarze nie niższym niż połowa pełnego etatu
[ ] urodzenia dziecka/dzieci po \Pole{data porodu}

DANE DZIECKA/DZIECI
\Pole{dziecko imie}, \Pole{data porodu}, \Pole{adres miejscowosc}

DANE DRUGIEGO RODZICA
\Pole{drugi rodzic}, PESEL \Pole{pesel}

ŚWIADCZENIA PROSZĘ PRZEKAZYWAĆ
[ ] w Polsce na rachunek bankowy
[ ] za granicę --- na rachunek bankowy (oprócz USA)

\vfill
{\footnotesize\color{szary}\itshape Zgodnie z art. 30 ustawy zasiłkowej wniosek o zasiłek za okres urlopu rodzicielskiego składa się w terminie do 21 dni po porodzie, jeśli chce się zasiłek za cały okres; późniejsze wnioski --- za okres od dnia złożenia. }

\clearpage

\section{Oświadczenie drugiego rodzica o niekorzystaniu z urlopu rodzicielskiego}

{\small\color{szary}Szablon P04 \textbullet{} art. 31 ust. 3a ustawy zasiłkowej}

\NaglowekPisma{ZAKŁAD UBEZPIECZEŃ SPOŁECZNYCH}

\NaglowekPisma{OŚWIADCZENIE}

Oświadczam, że jako ojciec dziecka \Pole{dziecko imie}, urodzonego dnia
\Pole{data porodu}, nie będę korzystał z przysługującego mi urlopu
rodzicielskiego w okresie, za jaki zasiłek macierzyński pobiera matka
dziecka --- \Pole{drugi rodzic}.

Jestem świadomy odpowiedzialności karnej za złożenie fałszywego oświadczenia.

\Pole{data pisma}
.....................................
(podpis)

\vfill
{\footnotesize\color{szary}\itshape Oświadczenie jest wymagane, gdy jeden z rodziców pobiera zasiłek za okres urlopu rodzicielskiego, a drugi z tego urlopu w tym czasie nie korzysta. }

\clearpage

\section{Wniosek o udzielenie urlopu ojcowskiego}

{\small\color{szary}Szablon P05 \textbullet{} art. 182\textsuperscript{3} \S{} 1 k.p.; art. 83 ustawy zasiłkowej}

\Pole{imie}\\[-2pt]
\Pole{adres}\\[-2pt]
PESEL: \Pole{pesel}

\Pole{platnik nazwa}

\Pole{data pisma}

\NaglowekPisma{WNIOSEK O UDZIELENIE URLOPU OJCOWSKIEGO}

Na podstawie art. 182\textsuperscript{3} \S{} 1 Kodeksu pracy wnoszę o udzielenie mi urlopu
ojcowskiego w wymiarze \Pole{tygodnie} tygodni, w okresie od \Pole{urlop od} do \Pole{urlop do}.

Urlopu ojcowskiego udziela się na pisemny wniosek ojca, składany w terminie
nie krótszym niż 7 dni przed rozpoczęciem urlopu.

Dziecko: \Pole{dziecko imie}, urodzone \Pole{data porodu}.

.....................................
(podpis pracownika)

\vfill
{\footnotesize\color{szary}\itshape Urlop ojcowski przysługuje w wymiarze 2 tygodni, nie dłużej niż do ukończenia przez dziecko 12 miesięcy (art. 182\textsuperscript{3} \S{} 1 i 4 k.p.). Można go dzielić na dwie części. Ojcowie adopcyjni i zastępczy mają odrębne zasady (art. 182\textsuperscript{4} k.p.). }

\clearpage

\section{Wniosek ojca o przejęcie części urlopu rodzicielskiego --- 9 tygodni}

{\small\color{szary}Szablon P06 \textbullet{} art. 182\textsuperscript{1}a \S{} 2 k.p.; art. 30 ust. 1 ustawy zasiłkowej}

\NaglowekPisma{WNIOSEK O UDZIELENIE URLOPU RODZICIELSKIEGO (CZĘŚĆ NIEPRZENOSZALNA)}

Na podstawie art. 182\textsuperscript{1}a \S{} 2 Kodeksu pracy wnoszę o udzielenie mi urlopu
rodzicielskiego w wymiarze \Pole{tygodnie} tygodni, w okresie od \Pole{urlop od} do \Pole{urlop do}.

Matka dziecka, \Pole{drugi rodzic}, wykorzystała \Pole{tygodnie matka} tygodni urlopu
macierzyńskiego, a z pozostałej części urlopu rodzicielskiego rezygnuje na
moją rzecz. Rezygnacja matki z urlopu rodzicielskiego na rzecz ojca wymaga:

[ ] pisemnego wniosku matki o rezygnację z urlopu rodzicielskiego w tej części,
[ ] mojego wniosku o udzielenie urlopu rodzicielskiego.

\begin{warunek}
Matka dziecka nie jest objęta ubezpieczeniem chorobowym, dlatego przysługuje\\[-2pt]
mi urlop rodzicielski na podstawie art. 182\textsuperscript{1}a \S{} 2\textsuperscript{1} Kodeksu pracy.
\end{warunek}

Dziecko: \Pole{dziecko imie}, urodzone \Pole{data porodu}.

\Pole{data pisma}
.....................................
(podpis pracownika)

\vfill
{\footnotesize\color{szary}\itshape Każdy z rodziców ma prawo do 9 tygodni urlopu rodzicielskiego, których nie można przenieść na drugiego rodzica (art. 182\textsuperscript{1}a \S{} 2 k.p.). Matka musi złożyć wniosek o rezygnację z odpowiedniej części urlopu (P07). }

\clearpage

\section{Wniosek matki o rezygnację z urlopu rodzicielskiego na rzecz ojca}

{\small\color{szary}Szablon P07 \textbullet{} art. 182\textsuperscript{1}a \S{} 2 k.p.; \S{} 4 rozporządzenia o wnioskach}

\NaglowekPisma{WNIOSEK O REZYGNACJĘ Z URLOPU RODZICIELSKIEGO}

Na podstawie art. 182\textsuperscript{1}a \S{} 2 Kodeksu pracy oświadczam, że rezygnuję z urlopu
rodzicielskiego w wymiarze \Pole{tygodnie} tygodni, w okresie od \Pole{urlop od} do \Pole{urlop do},
na rzecz ojca dziecka --- \Pole{drugi rodzic}.

\begin{warunek}
Macierzyński urlop został przez mnie wykorzystany w wymiarze co najmniej\\[-2pt]
14 tygodni po porodzie, a ojciec dziecka złożył wniosek o przyznanie mu\\[-2pt]
prawa do zasiłku macierzyńskiego za pozostały okres.
\end{warunek}

\begin{warunek}
Posiadam orzeczenie o niezdolności do samodzielnej egzystencji, dlatego\\[-2pt]
mogę zrezygnować z urlopu macierzyńskiego po wykorzystaniu 8 tygodni.
\end{warunek}

\Pole{data pisma}
.....................................
(podpis)

\vfill
{\footnotesize\color{szary}\itshape Matka może zrezygnować z urlopu macierzyńskiego po 14 tygodniach od porodu (art. 180 \S{} 4 k.p.), a w przypadku orzeczenia o niezdolności do samodzielnej egzystencji --- po 8 tygodniach (art. 180 \S{} 6 k.p.). }

\clearpage

\section{Wniosek o wypłatę zasiłku chorobowego w ciąży (Z-3)}

{\small\color{szary}Szablon P08 \textbullet{} art. 8, art. 11 ustawy zasiłkowej}

[FORMULARZ ZUS --- Z-3 / Z-3a / Z-3b]\\[-2pt]
Zaświadczenie płatnika składek

Pracownik występuje o:
[ ] zasiłek chorobowy            od \Pole{data zdarzenia} do \Pole{urlop do}
[ ] zasiłek opiekuńczy
[ ] zasiłek macierzyński
[ ] zasiłek w wysokości zasiłku macierzyńskiego
[ ] zasiłek wyrównawczy
[ ] świadczenie rehabilitacyjne

INFORMACJE O PRACOWNIKU
1. Zatrudniony od: \Pole{data zdarzenia}     w wymiarze czasu pracy: \Pole{wymiar etatu}
2. Czy zatrudniony na podstawie umowy o pracę tymczasową: TAK / NIE
3. Czy wykonuje pracę za granicą: TAK / NIE
4. Czy niezdolność do pracy powstała z powodu: choroby zawodowej / wypadku
przy pracy / wypadku w drodze do pracy lub z pracy / NIE
5. Czy w okresie niezdolności do pracy korzysta z urlopu bezpłatnego /
wychowawczego: TAK / NIE
6. Pracownikowi udzielono (jeśli występuje o zasiłek macierzyński):
[ ] urlopu macierzyńskiego      [ ] urlopu rodzicielskiego
[ ] uzupełniającego urlopu       [ ] urlopu na warunkach urlopu macierzyńskiego
[ ] urlopu ojcowskiego
od \Pole{urlop od} do \Pole{urlop do}
7. Czy pracownik w czasie urlopu rodzicielskiego wykonuje pracę: TAK / NIE

\vfill
{\footnotesize\color{szary}\itshape Zasiłek chorobowy w ciąży przysługuje w pierwszym dniu niezdolności do pracy w 100\% podstawy (art. 11 ust. 2 pkt 1 ustawy zasiłkowej). Formularz Z-3 wypełnia płatnik; zwolnienie lekarskie oznacza się kodem B (ciąża). }

\clearpage

\section{Wniosek o urlop macierzyński --- urodzenie martwego dziecka}

{\small\color{szary}Szablon P09 \textbullet{} art. 180 \S{} 4 pkt 3 k.p.; art. 29a ustawy zasiłkowej}

\NaglowekPisma{WNIOSEK O UDZIELENIE URLOPU MACIERZYŃSKIEGO}

Zgodnie z art. 180 \S{} 1 Kodeksu pracy zwracam się z wnioskiem o udzielenie mi
urlopu macierzyńskiego w wymiarze 8 tygodni od dnia porodu, przysługującego
w związku z urodzeniem martwego dziecka w dniu \Pole{data porodu}.

Urlop macierzyński będzie trwał od dnia \Pole{urlop od} do dnia \Pole{urlop do}.

Do wniosku załączam zaświadczenie szpitalne, że dziecko urodziło się martwe.

\Pole{data pisma}
.....................................
(podpis)

\vfill
{\footnotesize\color{szary}\itshape Przy urodzeniu martwego dziecka urlop macierzyński wynosi 8 tygodni (art. 180 \S{} 4 pkt 3 k.p.). Przysługuje też zasiłek pogrzebowy na dziecko martwo urodzone (art. 79 ust. 1 pkt 1 lit. b ustawy zasiłkowej). }

\clearpage

\section{Wniosek o urlop macierzyński po zgonie dziecka}

{\small\color{szary}Szablon P10 \textbullet{} art. 180 \S{} 4 pkt 3 k.p.; art. 29a ustawy zasiłkowej}

\NaglowekPisma{WNIOSEK O UDZIELENIE URLOPU MACIERZYŃSKIEGO}

Zgodnie z art. 180 \S{} 4 pkt 3 Kodeksu pracy wnoszę o udzielenie mi urlopu
macierzyńskiego w wymiarze \Pole{tygodnie} tygodni, przysługującego w związku
z urodzeniem dziecka, które zmarło w dniu \Pole{data zdarzenia}.

Urlop macierzyński ma trwać od dnia \Pole{urlop od} do dnia \Pole{urlop do}.

Do wniosku załączam skrócony odpis aktu urodzenia z adnotacją o zgonie dziecka.

\Pole{data pisma}
.....................................
(podpis)

\vfill
{\footnotesize\color{szary}\itshape Urlop wynosi 8 tygodni, nie krócej niż do 8 tygodni po porodzie.}

\clearpage

\section{Wniosek o zasiłek w wysokości zasiłku macierzyńskiego (likwidacja / upadłość)}

{\small\color{szary}Szablon P11 \textbullet{} art. 30 ust. 4, art. 31 ust. 1 ustawy zasiłkowej}

[FORMULARZ ZUS --- ZAM]\\[-2pt]
Wniosek o zasiłek w wysokości zasiłku macierzyńskiego

Wnoszę o wypłatę zasiłku w wysokości zasiłku macierzyńskiego za okres
urlopu macierzyńskiego / rodzicielskiego od \Pole{urlop od} do \Pole{urlop do}
w związku z niewypłaceniem przez pracodawcę świadczeń, których obowiązek
wypłaty przeszedł na Zakład Ubezpieczeń Społecznych z powodu:

[ ] ogłoszenia upadłości pracodawcy
[ ] likwidacji pracodawcy
[ ] śmierci pracodawcy

Pracodawca: \Pole{platnik nazwa}, NIP \Pole{platnik nip}.

\vfill
{\footnotesize\color{szary}\itshape Zasiłek wypłaca ZUS, gdy pracodawca nie wypłaca świadczeń z powodu upadłości, likwidacji lub śmierci pracodawcy (art. 30 ust. 4 ustawy zasiłkowej). }

\clearpage

\section{Wniosek o urlop na warunkach urlopu macierzyńskiego (przysposobienie)}

{\small\color{szary}Szablon P12 \textbullet{} art. 183 \S{} 1 k.p.; art. 29a, art. 30 ustawy zasiłkowej}

\NaglowekPisma{WNIOSEK O UDZIELENIE URLOPU NA WARUNKACH URLOPU MACIERZYŃSKIEGO}

Na podstawie art. 183 \S{} 1 Kodeksu pracy wnoszę o udzielenie mi urlopu na
warunkach urlopu macierzyńskiego w wymiarze \Pole{tygodnie} tygodni, w okresie
od \Pole{urlop od} do \Pole{urlop do}, w związku z:

[ ] przyjęciem dziecka na wychowanie w celu przysposobienia,
[ ] przyjęciem dziecka na wychowanie w ramach rodziny zastępczej
(innej niż zawodowa).

Dziecko: \Pole{dziecko imie}, urodzone \Pole{data porodu}.
Przyjęcie dziecka na wychowanie nastąpiło dnia \Pole{data zdarzenia}.

\begin{warunek}
Urlop przysługuje mi w wymiarze 20 tygodni (przy jednym dziecku), a przy\\[-2pt]
przyjęciu dziecka w wieku do 3 lat --- do ukończenia przez dziecko 4 lat.
\end{warunek}

\begin{warunek}
Urlop przysługuje mi w wymiarze wynikającym z art. 183 \S{} 1 k.p.;\\[-2pt]
rodzina zastępcza zawodowa nie ma prawa do tego urlopu.
\end{warunek}

\Pole{data pisma}
.....................................
(podpis)

\vfill
{\footnotesize\color{szary}\itshape Adopcja i rodzina zastępcza mają RÓŻNE wymiary urlopu. \quotedblbase{}Za życiem\textquotedblright{} (art. 183 \S{} 4 pkt 3 i \S{} 4\textsuperscript{1} k.p.) dotyczy rodziny zastępczej, nie przysposobienia. }

\clearpage

\section{Wniosek o zasiłek opiekuńczy (Z-15A)}

{\small\color{szary}Szablon P13 \textbullet{} art. 32--35 ustawy zasiłkowej}

[FORMULARZ ZUS --- Z-15A / Z-15B]\\[-2pt]
Wniosek o zasiłek opiekuńczy

Wnoszę o wypłatę zasiłku opiekuńczego z powodu opieki nad:

[ ] dzieckiem zdrowym w wieku do 8 lat (art. 32 ustawy zasiłkowej)
[ ] chorym dzieckiem w wieku do 14 lat
[ ] chorym dzieckiem w wieku do 18 lat (z orzeczeniem o niepełnosprawności)
[ ] innym chorym członkiem rodziny

Okres opieki: od \Pole{urlop od} do \Pole{urlop do}.

Zasiłek opiekuńczy przysługuje w wymiarze:
- 60 dni w roku kalendarzowym (dzieci),
- 14 dni w roku kalendarzowym (inni członkowie rodziny).

\vfill
{\footnotesize\color{szary}\itshape Zasiłek opiekuńczy za opiekę nad dzieckiem: 80\% podstawy. Nie wlicza się do okresu zasiłkowego 182 dni (art. 34 ustawy zasiłkowej). }

\clearpage

\section{Wniosek o łączenie urlopu rodzicielskiego z pracą}

{\small\color{szary}Szablon P14 \textbullet{} art. 182\textsuperscript{1}f k.p.; art. 31 ust. 3a ustawy zasiłkowej}

\NaglowekPisma{WNIOSEK O WYRAŻENIE ZGODY NA ŁĄCZENIE URLOPU RODZICIELSKIEGO Z PRACĄ}

Na podstawie art. 182\textsuperscript{1}f Kodeksu pracy wnoszę o wyrażenie zgody na wykonywanie
pracy w czasie urlopu rodzicielskiego, w okresie od \Pole{urlop od} do \Pole{urlop do},
w wymiarze \Pole{wymiar etatu} etatu.

Łączenie urlopu rodzicielskiego z pracą jest możliwe przez cały okres
przysługującego urlopu, z tym że praca może być wykonywana w wymiarze
nie wyższym niż 1/2 pełnego wymiaru (a w szczególnych przypadkach --- wyższym).

Pracodawca może odmówić wyrażenia zgody, jeżeli praca w niepełnym wymiarze
nie jest możliwa; odmowa wymaga uzasadnienia.

\Pole{data pisma}
.....................................
(podpis)

\vfill
{\footnotesize\color{szary}\itshape Przy łączeniu urlopu z pracą zasiłek macierzyński ulega proporcjonalnemu zmniejszeniu. Przepis obowiązuje od 1 lipca 2026 r. }

\clearpage

\section{Wniosek o zmianę wymiaru urlopu rodzicielskiego (korekta łączenia)}

{\small\color{szary}Szablon P15 \textbullet{} art. 182\textsuperscript{1}b k.p.; art. 30, art. 32 ustawy zasiłkowej}

\NaglowekPisma{WNIOSEK O ZMIANĘ WYMIARU WYKORZYSTANIA URLOPU RODZICIELSKIEGO}

Na podstawie \S{} 4 pkt 1 lit. b) i c) rozporządzenia MRPiPS w sprawie wniosków
dotyczących uprawnień pracowników związanych z rodzicielstwem wnoszę o zmianę
wymiaru wykonywania pracy w trakcie urlopu rodzicielskiego:

[ ] z \Pole{tygodnie matka} tygodni na \Pole{tygodnie} tygodni od \Pole{urlop od}
[ ] z \Pole{tygodnie matka} tygodni na \Pole{tygodnie} tygodni od \Pole{urlop do}

\Pole{data pisma}
.....................................
(podpis)

\vfill
{\footnotesize\color{szary}\itshape Zmianę wymiaru urlopu (bez zmiany jego łącznego wymiaru) dopuszcza art. 182\textsuperscript{1}b \S{} 1 k.p.}

\clearpage

\section{Wniosek o uzupełniający urlop macierzyński (wcześniak / szpital)}

{\small\color{szary}Szablon P16 \textbullet{} art. 180\textsuperscript{1} k.p.; art. 29 ust. 5 ustawy zasiłkowej}

\NaglowekPisma{WNIOSEK O UDZIELENIE UZUPEŁNIAJĄCEGO URLOPU MACIERZYŃSKIEGO}

Na podstawie art. 183\textsuperscript{1} \S{} 1 pkt 1 Kodeksu pracy wnoszę o udzielenie mi
uzupełniającego urlopu macierzyńskiego w wymiarze stanowiącym różnicę
między 41 tygodniami a wymiarami urlopów wykorzystanych, w okresie
od \Pole{urlop od} do \Pole{urlop do}.

Uzupełniający urlop macierzyński przysługuje w wymiarze:

dzień porodu do 28. tygodnia ciąży        = 41 tygodni − wykorzystane tygodnie
od 28. do 36. tygodnia ciąży             = 41 tygodni − wykorzystane tygodnie

Tydzień ciąży w chwili porodu: \Pole{tydzien ciazy}
Masa urodzeniowa dziecka: \Pole{masa} g
Pobyt dziecka w szpitalu: od \Pole{szpital od} do \Pole{szpital do}  (\Pole{dni hospitalizacji} dni)

\Pole{data pisma}
.....................................
(podpis)

\vfill
{\footnotesize\color{szary}\itshape Uzupełniający urlop macierzyński przysługuje w razie urodzenia wcześniaka albo gdy dziecko wymagało hospitalizacji po porodzie (art. 183\textsuperscript{1} \S{} 1 pkt 1--2 k.p.). }

\clearpage

\end{document}
